\subsection{Summary of findings}
\label{sec:summary}

\begin{enumerate}
    \item \textbf{Text read from scratch fails across projects.} AUC is
          \mBvGgnnLexAuc{} (BigVul) and \mDvGgnnLexAuc{} (DiverseVul), close to
          random.
    \item \textbf{Removing project vocabulary helps by about \numdeltaMin{} to \numdeltaMax{} AUC,}
          on both corpora and under every architecture we ran (the Transformer on
          BigVul only), whether it is removed
          by anonymising identifiers or by using the grammar kind.
    \item \textbf{The grammar kind does not beat anonymised text;} its advantage
          is that it needs no preprocessing decision.
    \item \textbf{The engineered operation taxonomy adds nothing} to the grammar
          kind, and the reason is measurable: it is almost a function of it
          (Section~\ref{sec:rq_why}). As a static analysis it is accurate for the
          question it answers (memory operation or not: precision
          \numpooledMemopP{}, recall \numpooledMemopR{}) but less accurate than
          first reported at the level of individual classes.
    \item \textbf{Much of that gain is about function size.} Within bins of
          function size the gain over text is not distinguishable from zero for
          the grammar kind in every corpus--architecture cell we ran, and within-bin AUC is \mBvGgnnKindSs{} (BigVul) and
          \mDvGgnnKindSs{} (DiverseVul).
    \item \textbf{Classifiers that read no code beat every network:} a
          one-line size rule on both corpora and, on BigVul, a classifier on
          formatting statistics alone (AUC \rBvFormatAuc{}), because the corpus
          has a formatting artefact. Models that read raw text partly exploit
          it (Section~\ref{sec:artefact}). Absolute performance of the graph
          models is weak.
\end{enumerate}
